\section{Công trình liên quan}
\label{sec:lien-quan}

\paragraph{Các chiều chất lượng dữ liệu.} Chất lượng dữ liệu từ lâu được mô tả theo nhiều chiều như độ đầy đủ, độ chính xác, tính nhất quán và tính kịp thời \citep{wang1996beyond,batini2009methodologies}. Các khung đánh giá truyền thống đo từng chiều độc lập và để người dùng tự quyết ngưỡng chấp nhận. Chúng không trả lời câu hỏi một mức lỗi cụ thể làm mô hình học máy tệ đi bao nhiêu.

\paragraph{\citet{mohammed2025effects}: ảnh hưởng của chất lượng dữ liệu là theo tác vụ.} Đây là bài gốc chính của đề tài. Nhóm tác giả định nghĩa sáu chiều chất lượng cho dữ liệu dạng bảng, mỗi chiều có một bộ ``polluter'' làm bẩn dữ liệu theo mức độ điều khiển được. Họ chạy 19 thuật toán trên nhiều bộ dữ liệu cho ba tác vụ phân loại, hồi quy và phân cụm. Kết luận chính: cùng một loại lỗi, cùng một mức độ, nhưng hiệu năng giảm rất khác nhau giữa các tác vụ và thuật toán. Bài báo dừng ở đường cong quan sát, không có mô hình dự đoán. Các bộ dữ liệu cũng nhỏ (khoảng 1.000 đến 25.000 dòng), chưa đặt ra bài toán quy mô. Đề tài kế thừa định nghĩa các chiều chất lượng, cách làm bẩn và kịch bản ``làm bẩn cả train lẫn test'' của bài này.

\paragraph{\citet{bayram2023dqsops}: DQSOps, dự đoán nhanh một điểm chất lượng.} DQSOps đề xuất quy trình chấm điểm chất lượng dữ liệu liên tục trong hệ thống vận hành. Các chiều chất lượng được chuẩn hóa rồi gộp bằng PCA thành một điểm duy nhất; một mô hình học máy sau đó học cách dự đoán nhanh điểm này để khỏi phải tính lại toàn bộ. Điểm mạnh là có mô hình dự đoán và hướng đến vận hành. Tuy vậy, thực nghiệm chỉ gồm một use case công nghiệp với dữ liệu chuỗi thời gian đơn biến. Điểm yếu chính, cũng là chỗ đề tài nhắm vào: nhãn huấn luyện là chính điểm PCA tự định nghĩa, không gắn với hiệu năng mô hình phía sau, và dùng chung cho mọi tác vụ. Đề tài cài lại cách tính điểm này làm baseline quan trọng nhất (Mục~\ref{sec:baseline}).

\paragraph{Profiling dữ liệu quy mô lớn.} Khi dữ liệu vượt bộ nhớ một máy, các thư viện xử lý trong bộ nhớ như pandas \citep{mckinney2010pandas} không còn dùng được. Apache Spark \citep{zaharia2016spark} chia dữ liệu thành phân vùng và xử lý dần nên chạy được trong giới hạn bộ nhớ cố định, đổi lại tốn chi phí khởi động. Đề tài dùng Spark để kiểm chứng rằng bước profiling, bước duy nhất phải chạm vào toàn bộ dữ liệu khi dùng khung đánh giá, mở rộng được.

\paragraph{Vị trí của đề tài.} Đề tài ghép hai hướng trên: lấy cách đo và cách làm bẩn có kiểm soát của \citet{mohammed2025effects} để sinh dữ liệu huấn luyện, rồi xây một mô hình dự đoán như DQSOps, nhưng nhãn là mức suy giảm hiệu năng thật và mô hình biết loại tác vụ.
